\section{The involution saddle at every fixed order}\label{sec:involutions}
This section records a classical saddle calculation in a form directly usable with Theorem~\ref{thm:main}. The method and all-order involution asymptotics are prior work \cite{MW1955,MW1956}. The details below supply a controlled remainder and an explicit finite coefficient rule.

Since $I_n=n![z^n]\exp(z+z^2/2)$, the positive saddle is
\begin{equation}\label{eq:saddle}
r=\frac{\sqrt{4n+1}-1}{2},\qquad n=r^2+r.
\end{equation}
\begin{proposition}\label{prop:saddle}
For every fixed integer $J\geq0$,
\begin{equation}\label{eq:saddleexpansion}
I_n=\frac{n!e^{r+r^2/2}}{2\sqrt\pi\,r^{n+1}}
 \left(\sum_{j=0}^{J}\gamma_jr^{-j}+O_J(r^{-J-1})\right).
\end{equation}
The rational coefficients are given by the finite recursion
\begin{align}
p_m(v)&=\frac{(\ii v)^{m+1}}{(m+1)!}
            +\frac{2^{m+1}(\ii v)^{m+2}}{(m+2)!},\quad m\geq1,\notag\\
Q_0(v)&=1,\qquad
Q_j(v)=\frac1j\sum_{m=1}^j mp_m(v)Q_{j-m}(v),\label{eq:saddlerec}\\
\gamma_j&=\frac1{\sqrt\pi}\int_{\R}e^{-v^2}Q_j(v)\,\dd v.\notag
\end{align}
The integral is an algebraic rule: send odd powers of $v$ to zero and $v^{2k}$ to $(2k-1)!!/2^k$.
\end{proposition}
\begin{proof}
Cauchy's formula on $|z|=r$ gives the prefactor $n!e^{r+r^2/2}/(2\pi r^n)$ times the integral of $e^{\Phi_r(\theta)}$ on $[-\pi,\pi]$, where
\[
\Phi_r(\theta)=r(e^{\ii\theta}-1-\ii\theta)
 +\tfrac12r^2(e^{2\ii\theta}-1-2\ii\theta).
\]
Its real part is exactly
\begin{equation}\label{eq:realphase}
\Re\Phi_r(\theta)=-r(1-\cos\theta)-r^2\sin^2\theta.
\end{equation}
For $|\theta|\leq\pi/2$ this is at most $-4r^2\theta^2/\pi^2$; for $|\theta|\geq\pi/2$ it is at most $-r$. Fix $0<\eps<1/3$. The integral outside $|\theta|\leq r^{-1+\eps}$ is smaller than every fixed algebraic order. The second inequality also controls the possible secondary peak near $\theta=\pi$.

On the central interval put $\theta=v/r$ and $h=1/r$. Direct Taylor expansion gives
\[
\Phi_r(v/r)=-v^2+\sum_{m\geq1}h^mp_m(v).
\]
The coefficient polynomials of the exponential satisfy \eqref{eq:saddlerec}. For fixed $J$, Taylor's theorem in $h$ bounds the remainder on $|v|\leq r^\eps$ by $h^{J+1}$ times a fixed polynomial in $|v|$ times $e^{\delta v^2}$, with any chosen $\delta>0$ once $r$ is large. Indeed, derivatives of the analytic exponent of fixed order are bounded by fixed polynomials in $|v|$ when $h|v|=o(1)$, and its non-Gaussian part is bounded by $Ch(v^2+|v|^3)\leq\delta v^2$. Taking $\delta<1$ leaves an integrable Gaussian majorant after multiplication by $e^{-v^2}$. Extending polynomial Gaussian integrals to all of $\R$ adds exponentially small error. Finally $\dd\theta=\dd v/r$ gives \eqref{eq:saddleexpansion}.
\end{proof}

The first saddle coefficients are
\begin{equation}\label{eq:gammas}
(\gamma_0,\ldots,\gamma_6)=
\left(1,-\frac14,-\frac7{96},\frac{77}{384},
-\frac{3989}{18432},\frac{3275}{24576},\frac{1076179}{26542080}\right).
\end{equation}

\subsection{Conversion to half integer powers}
Write $t=n^{-1/2}$ and
\begin{equation}\label{eq:Bofn}
B(n)=\frac1{\sqrt2}\left(\frac ne\right)^{n/2}e^{\sqrt n-1/4}.
\end{equation}
For every fixed $J$,
\begin{equation}\label{eq:Ialpha}
I_n=B(n)\left(\sum_{j=0}^{J}\alpha_jn^{-j/2}
                  +O_J(n^{-(J+1)/2})\right),
\end{equation}
where
\begin{align}
\alpha_0&=1,&\alpha_1&=\frac7{24},&
\alpha_2&=-\frac{119}{1152},&
\alpha_3&=-\frac{7933}{414720},\notag\\
\alpha_4&=\frac{1967381}{39813120},&
\alpha_5&=-\frac{57200419}{1337720832},&
\alpha_6&=\frac{6340449533}{687970713600}.\label{eq:alphas}
\end{align}
The first three nonconstant coefficients also appear explicitly in Bornemann \cite[Appendix C, (120)]{Bornemann}. The logarithmic form is
\begin{align}
\log I_n={}&\frac n2(\log n-1)+\sqrt n-\frac14-\frac12\log2\notag\\
 &+\frac7{24}t-\frac7{48}t^2+\frac{37}{1920}t^3
   +\frac5{128}t^4-\frac{5879}{107520}t^5
   +\frac{43}{1440}t^6+O(t^7).\label{eq:logI}
\end{align}

For an all-order rule, put
$R(t)=\sqrt{1+t^2/4}-t/2$, so that $r=R(t)/t$ and $\log R(t)=-\operatorname{arsinh}(t/2)$. Let $\mathcal B(t)$ denote the formal logarithmic correction to $\log B(n)$. Substitution in \eqref{eq:saddleexpansion} and Stirling's expansion give
\begin{align}
\mathcal B(t)={}&\frac{R(t)}{2t}
 +\frac{\operatorname{arsinh}(t/2)}{t^2}
 +\operatorname{arsinh}(t/2)-\frac1t+\frac14\notag\\
 &+\sum_{k\geq1}\frac{B_{2k}}{2k(2k-1)}t^{4k-2}
 +\log\left(\sum_{j\geq0}\gamma_j\left(\frac{t}{R(t)}\right)^j\right).
 \label{eq:allorderI}
\end{align}
Here $B_{2k}$ are Bernoulli numbers; the apparent negative powers cancel. At any specified order, every sum and series operation is truncated, so the rule is finite. Then $\alpha_j=[t^j]\exp\mathcal B(t)$. Proposition~\ref{prop:saddle} and Stirling's remainder justify each fixed truncation.

\subsection{Ratios at fixed shifts}
For any fixed nonnegative integer $s$, subtraction of sufficiently long logarithmic expansions gives
\begin{align}
\frac{I_{n-s}}{I_n}=n^{-s/2}\bigg[1
 &-\frac s2t+\frac{3s^2}{8}t^2
 +\frac{-7s^3-6s^2+7s}{48}t^3\notag\\
 &+\frac{25s^4+56s^3-28s^2-56s}{384}t^4\notag\\
 &+\frac{-81s^5-340s^4-30s^3+700s^2+111s}{3840}t^5
 +O_s(t^6)\bigg].\label{eq:shiftratio}
\end{align}
For all orders, exponentiate the finite truncation of
\begin{align}
\mathcal D_s(t)={}&\frac{t^{-2}-s}{2}\log(1-st^2)+\frac s2
 +\frac{\sqrt{1-st^2}-1}{t}\notag\\
 &+\mathcal B\left(\frac{t}{\sqrt{1-st^2}}\right)-\mathcal B(t).
 \label{eq:shiftrule}
\end{align}
Then $I_{n-s}/I_n$ has asymptotic series $n^{-s/2}\exp\mathcal D_s(t)$. Uniformity is asserted only for $s$ in a prescribed finite set. This argument retains enough terms before subtracting; it does not differentiate an unproved remainder.
